% GENERATED FILE. Edit canonical lessons, then run npm run generate:books.
% canonical-source-hash: fnv1a64:e070eacfe04522ef
% canonical-lessons: IT-C483-livido, IT-C483-graffio, IT-C483-scottatura, IT-C483-inciampare, IT-C483-tagliarsi

\chapter{Piccoli incidenti}
\label{ch:it-piccoli-incidenti}
\hlchaptermodality{483}

\begin{chapteropening}
\textbf{By the end of this chapter:} \emph{I can talk about small injuries: bruises, scratches, burns and cuts.}

Complete the last lesson of chapter 483: I can talk about small injuries: bruises, scratches, burns and cuts.
\end{chapteropening}

\section[il livido]{il livido — bruise}

\label{lesson:IT-C483-livido}



\begin{quote}
Before the new one: say the Italian for temple, then the Italian for jaw.
\end{quote}

\subsection*{You'll want to know: il livido}
\textbf{il livido} — \textquotedblleft{}bruise\textquotedblright{}.

\begin{quote}
Sono caduto e ho un livido sul ginocchio.
\end{quote}

\begin{quote}
\emph{I fell and I have a bruise on my knee.}
\end{quote}

Also an adjective: \emph{livido} — bluish, or pale with anger.

\begin{cousinweb}
\emph{livido} is a learned borrowing from Latin \emph{līvidus} \textquotedblleft{}bluish, black and blue\textquotedblright{} (EtymDB; Lewis \& Short), from \emph{līvēre} \textquotedblleft{}to be bluish\textquotedblright{}.
\end{cousinweb}

\begin{grammarlens}[title={What you've built}]
One more word for telling, explaining and describing at B1.
\end{grammarlens}

\subsection*{Your turn}
\begin{itemize}
\raggedright
  \item \emph{Say it:} \emph{il livido}
  \item \emph{Say it:} \emph{il livido}, once more
  \item \emph{Say it:} a sentence of your own with \emph{il livido}
  \item \emph{Recall:} say the Italian for temple, then the Italian for jaw, then say \emph{il livido} again
\end{itemize}

\begin{tcolorbox}[breakable,colback=teal!4,colframe=teal!35!black,title={Before you move on}]
What does il livido mean? (\textquotedblleft{}Bruise\textquotedblright{}.) Say it once more.
\end{tcolorbox}
\section[il graffio]{il graffio — scratch}

\label{lesson:IT-C483-graffio}



\begin{quote}
Before the new one: say the Italian for jaw, then the Italian for bruise.
\end{quote}

\subsection*{You'll want to know: il graffio}
\textbf{il graffio} — \textquotedblleft{}scratch\textquotedblright{}.

\begin{quote}
Il gatto mi ha fatto un graffio sulla mano.
\end{quote}

\begin{quote}
\emph{The cat scratched my hand.}
\end{quote}

The verb is \emph{graffiare} — to scratch.

\begin{cousinweb}
The origin of \emph{graffio} is not settled; it is often linked to a Germanic word for a hook or claw.
\end{cousinweb}

\begin{grammarlens}[title={What you've built}]
One more word for telling, explaining and describing at B1.
\end{grammarlens}

\subsection*{Your turn}
\begin{itemize}
\raggedright
  \item \emph{Say it:} \emph{il graffio}
  \item \emph{Say it:} \emph{il graffio}, once more
  \item \emph{Say it:} a sentence of your own with \emph{il graffio}
  \item \emph{Recall:} say the Italian for jaw, then the Italian for bruise, then say \emph{il graffio} again
\end{itemize}

\begin{tcolorbox}[breakable,colback=teal!4,colframe=teal!35!black,title={Before you move on}]
What does il graffio mean? (\textquotedblleft{}Scratch\textquotedblright{}.) Say it once more.
\end{tcolorbox}
\section[la scottatura]{la scottatura — burn; sunburn}

\label{lesson:IT-C483-scottatura}



\begin{quote}
Before the new one: say the Italian for bruise, then the Italian for scratch.
\end{quote}

\subsection*{You'll want to know: la scottatura}
\textbf{la scottatura} — \textquotedblleft{}burn; sunburn\textquotedblright{}.

\begin{quote}
Sono stato al sole troppo a lungo e ora ho una scottatura.
\end{quote}

\begin{quote}
\emph{I stayed in the sun too long and now I have sunburn.}
\end{quote}

From \emph{scottarsi} — to get burnt.

\begin{cousinweb}
\emph{scottatura} is built on \emph{scottare} \textquotedblleft{}to scald\textquotedblright{}, usually traced to Latin \emph{excoctāre}, from \emph{coquere} \textquotedblleft{}to cook\textquotedblright{} (Lewis \& Short).
\end{cousinweb}

\begin{grammarlens}[title={What you've built}]
One more word for telling, explaining and describing at B1.
\end{grammarlens}

\subsection*{Your turn}
\begin{itemize}
\raggedright
  \item \emph{Say it:} \emph{la scottatura}
  \item \emph{Say it:} \emph{la scottatura}, once more
  \item \emph{Say it:} a sentence of your own with \emph{la scottatura}
  \item \emph{Recall:} say the Italian for bruise, then the Italian for scratch, then say \emph{la scottatura} again
\end{itemize}

\begin{tcolorbox}[breakable,colback=teal!4,colframe=teal!35!black,title={Before you move on}]
What does la scottatura mean? (\textquotedblleft{}Burn; sunburn\textquotedblright{}.) Say it once more.
\end{tcolorbox}
\section[inciampare]{inciampare — to trip, to stumble}

\label{lesson:IT-C483-inciampare}



\begin{quote}
Before the new one: say the Italian for scratch, then the Italian for burn.
\end{quote}

\subsection*{You'll want to know: inciampare}
\textbf{inciampare} — \textquotedblleft{}to trip, to stumble\textquotedblright{}.

\begin{quote}
Sono inciampato sulle scale e sono caduto.
\end{quote}

\begin{quote}
\emph{I tripped on the stairs and fell.}
\end{quote}

It takes \emph{essere}: \emph{sono inciampato}. \emph{Inciampare in} something.

\begin{cousinweb}
The origin of \emph{inciampare} is not settled; it may be built on an old dialect word \emph{ciampa} \textquotedblleft{}paw, foot\textquotedblright{}.
\end{cousinweb}

\begin{grammarlens}[title={What you've built}]
One more word for telling, explaining and describing at B1.
\end{grammarlens}

\subsection*{Your turn}
\begin{itemize}
\raggedright
  \item \emph{Say it:} \emph{inciampare}
  \item \emph{Say it:} \emph{inciampare}, once more
  \item \emph{Say it:} a sentence of your own with \emph{inciampare}
  \item \emph{Recall:} say the Italian for scratch, then the Italian for burn, then say \emph{inciampare} again
\end{itemize}

\begin{tcolorbox}[breakable,colback=teal!4,colframe=teal!35!black,title={Before you move on}]
What does inciampare mean? (\textquotedblleft{}To trip, to stumble\textquotedblright{}.) Say it once more.
\end{tcolorbox}
\section[tagliarsi]{tagliarsi — to cut oneself}

\label{lesson:IT-C483-tagliarsi}



\begin{quote}
Before the new one: say the Italian for burn, then the Italian for to trip.
\end{quote}

\subsection*{You'll want to know: tagliarsi}
\textbf{tagliarsi} — \textquotedblleft{}to cut oneself\textquotedblright{}.

\begin{quote}
Mi sono tagliato con il coltello.
\end{quote}

\begin{quote}
\emph{I cut myself with the knife.}
\end{quote}

\emph{Tagliarsi un dito} — to cut a finger.

\begin{cousinweb}
\emph{tagliarsi} is the reflexive of \emph{tagliare} \textquotedblleft{}to cut\textquotedblright{}, inherited from Late Latin \emph{tāleāre}, from \emph{tālea} \textquotedblleft{}a rod, a cutting\textquotedblright{} (Lewis \& Short).
\end{cousinweb}

\begin{grammarlens}[title={What you've built}]
One more word for telling, explaining and describing at B1.
\end{grammarlens}

\subsection*{Your turn}
\begin{itemize}
\raggedright
  \item \emph{Say it:} \emph{tagliarsi}
  \item \emph{Say it:} \emph{tagliarsi}, once more
  \item \emph{Say it:} a sentence of your own with \emph{tagliarsi}
  \item \emph{Recall:} say the Italian for burn, then the Italian for to trip, then say \emph{tagliarsi} again
\end{itemize}

\begin{tcolorbox}[breakable,colback=teal!4,colframe=teal!35!black,title={Before you move on}]
What does tagliarsi mean? (\textquotedblleft{}To cut oneself\textquotedblright{}.) Say it once more.
\end{tcolorbox}
